\subsection{赋值的比较}

设 K 是域，A 是 K 的一个赋值环。对 K 中含 A 的每个子环 B，有 \(U(A) \subset U(B)\) 。于是存在一个典范同态 \(\lambda\) ，它由 \(\Gamma_{A} = K^{*}/U(A)\) 映到 \(\Gamma_{B} = K^{*}/U(B)\) 上，其核为 \(U(B)/U(A)\) 。设 \(v_{A}\) 与 \(v_{B}\) 分别表示由 A 与 B 定义的 K 上的典范赋值（ \(\S 3\) 第 2 节），则

\[
v _ {\mathrm{B}} = \lambda \circ v _ {\mathrm{A}}.\tag{1}
\]

由于 \(A \subset B\) ，\(\lambda\) 把 \(\Gamma_{A}\) 的正元素映到 \(\Gamma_{B}\) 的正元素，因此是递增的。于是（命题 3）\(H_{B}\) 作为 \(\lambda\) 的核，是 \(\Gamma_{A}\) 的孤立子群，并且 \(\lambda\) 分解为 \(\Gamma_{A} \to \Gamma_{A}/H_{B} \xrightarrow{\mu} \Gamma_{B}\) ，其中 \(\mu\) 是递增的双射同态，因而是全序群的同构；因此 \(\Gamma_{B}\) 等同于商全序群 \(\Gamma_{A}/H_{B}\) 。

命题 4. 映射 \(\mathbf{B} \mapsto \mathbf{H}_{\mathbf{B}}\) 是由 \(\mathbf{K}\) 中含 \(\mathbf{A}\) 的子环集合到 \(\Gamma_{\mathbf{A}}\) 的孤立子群集合上的一个递增双射。

给定 \(H_{B}, v_{B}\) 在等价的意义下就被确定，因而 B 被唯一确定。另一方面，设 H 是 \(\Gamma_{A}\) 的一个孤立子群；把 \(\Gamma_{A}/H\) 视为全序群（命题 3），复合映射

\[
\mathrm{K} ^ {*} \xrightarrow {v _ {\mathrm{A}}} \Gamma_ {\mathrm{A}} \longrightarrow \Gamma_ {\mathrm{A}} / \mathrm{H}
\]

是 K 上的一个赋值，其环含 A。

注. 在上述假设下，设 \(f\) 表示 \(\mathbf{B}\) 到 \(\kappa(\mathbf{B})\) 上的典范同态，\(\mathbf{A}' = f(\mathbf{A})\) ；它是 \(\kappa(\mathbf{B})\) 的一个赋值环（第 1 节命题 2）。于是 \(\overline{f}(\kappa(\mathbf{B})^*) = \mathbf{U}(\mathbf{B}), \overline{f}(\mathbf{A}') = \mathbf{A}, \overline{f}(\mathfrak{m}(\mathbf{A}')) = \mathfrak{m}(\mathbf{A})\) ，因而

\[
\overline {{f}} (\mathbf {U} (\mathbf {A} ^ {\prime})) = \mathbf {U} (\mathbf {A}).
\]

于是存在由 \(\mathbf{U}(\mathbf{B}) / \mathbf{U}(\mathbf{A})\) 到 \(\kappa (\mathbf{B})^{*} / \mathbf{U}(\mathbf{A}^{\prime}) = \Gamma_{\mathbf{A}}\) 上的典范同构。正合列

\[
0 \rightarrow \mathrm{U(B)} / \mathrm{U(A)} \rightarrow \Gamma_ {\mathrm{A}} \rightarrow \Gamma_ {\mathrm{B}} \rightarrow 0
\]

进而给出正合列

\[
0 \rightarrow \Gamma_ {A ^ {\prime}} \rightarrow \Gamma_ {A} \rightarrow \Gamma_ {B} \rightarrow 0.
\]

例. 设 \(k\) 是域，

\[
\mathrm{E} = k (\mathrm{X}) \quad \text { and } \quad \mathrm{K} = k (\mathrm{X}, \mathrm{Y}) = \mathrm{E} (\mathrm{Y})
\]

（X, Y 是未定元）。设 \(B = E[Y]_{(Y)}\) 是由主理想整环 E{[}Y{]} 的极端元 Y 所定义的 K 的赋值环（§ 1 第 4 节命题 3）。剩余域 \(\kappa(\mathbf{B})\) 典范地等同于 \(E[Y]/(Y) = E\) 。类似地，设 \(A' = k[X]_{(X)}\) 是 \(E = k(X)\) 的、由 \(k[X]\) 的极端元 X 所定义的赋值环。用 \(h_{B}\) 表示与 B 相伴的 E 的位，并记 \(A = \overline{h_{B}}(A')\) ，就定义了 K 的一个含于 B 的赋值环 A，且 \(\kappa(A) = \kappa(A') = k\) 。典范位 \(h_{A}: K \to k\) 可描述如下：若 \(f(X, Y)\) 是 K 的元素，则先在 f 中令 Y = 0（得到 \(\hat{E} = k(X)^{-}\) 的一个元素），再在所得结果中令 X = 0。群 \(\Gamma_{A'}\) 与 \(\Gamma_{B}\) 典范地同构于 Z（§ 3 第 4 节例 4）。*不难证明（参见~§ 10 第 2 节引理 2），群 \(\Gamma_{A}\) 同构于字典积 \(Z \times Z\)，且赋值 \(v_{A}\) 等价于 § 3 第 4 节例 6 末尾所定义的赋值。*

